\section{Introduction}
\label{sec:tpc38-introduction}

The orbit--sliced program in TPC--34--TPC--37 reduces a scalar prime--pair
remainder to a coefficient--specific energy estimate for a reflected
M\"obius divisor sum \citep{WangTPC34,WangTPC35,WangTPC36,WangTPC37}.
The physical row mask is now known to have subpolynomial signed projective
complexity, the zero residue faces have been removed at the target scale,
and all seven proper faces of a three--modulus completion close on the
exact equality.  The full face remains unsuppressed there.

TPC--37 recorded two possible continuations.  The first works after
finite--field completion and asks for an actual weighted Mellin--overlap
estimate, including its principal boundary.  The second keeps the integer
determinant
\begin{equation}
 F=su-N
 \label{eq:tpc38-intro-determinant}
\end{equation}
visible and asks for a target--coupled shifted--fiber estimate.  This paper
audits the second continuation.  The question is not whether one may write
another formal Hilbert--space norm.  It is whether the moving target and
the CRT faces create genuine orthogonality at the physical normalization.

\subsection{What the target coupling really is}

Let \(q_1,q_2,q_3\) be distinct odd primes, put
\begin{equation}
 M=q_1q_2q_3,\qquad
 \Phi=\varphi(M)=\prod_{i=1}^3(q_i-1),
 \label{eq:tpc38-intro-M-Phi}
\end{equation}
and let \(N\) be a unit modulo \(M\).  TPC--37 defines, in each local
coordinate, a principal face \(m_{i,N}(F)\) and a nonprincipal aggregate
\(c_{i,N}(F)\).  Their tensor products give eight faces.  The full face
\(C_N\) uses three nonprincipal coordinates; the proper union \(P_N\)
contains the other seven.

The first result of this paper is the exact identity
\begin{equation}
 K_{S,N}(F)
 =\frac1\Phi\sum_{\chi\in\mathfrak X_S}
 \chi(N+F)\overline{\chi(N)},
 \label{eq:tpc38-intro-character-face}
\end{equation}
where the local component of \(\chi\) is nonprincipal exactly for
\(i\in S\).  Since \(N+F=su\), the phase factors further as
\begin{equation}
 \chi(N+F)\overline{\chi(N)}
 =\chi(s)\chi(u)\overline{\chi(N)}.
 \label{eq:tpc38-intro-product-factorization}
\end{equation}
Zero--extended characters make this valid without a hidden correction
when \(s\) or \(u\) is nonunit; then every face vanishes in a killed local
coordinate.  On the regular support
\(N+F\in(\Z/M\Z)^\times\), the finite CRT target dependence is
diagonalized exactly by the unit--group ratio
\begin{equation}
 r=\frac{N+F}{N}\pmod M.
 \label{eq:tpc38-intro-ratio}
\end{equation}

For fixed \(F\), on its regular domain this ratio map is injective in every local coordinate in
which \(q_i\nmid F\), and is constant in every coordinate in which
\(q_i\mid F\).  Its exact largest fiber is
\begin{equation}
 d_M(F)=\prod_{q_i\mid F}(q_i-1).
 \label{eq:tpc38-intro-collision-multiplicity}
\end{equation}
When \(q_i\asymp J\), this gives the target collision scale ladder
\begin{equation}
 1,\qquad J,\qquad J^2,\qquad J^3,
 \label{eq:tpc38-intro-collision-ladder}
\end{equation}
according as \(F\) is divisible by zero, one, two, or three orbit--scale
primes.  The underlying multiplicity formula is exact; the displayed
\(J\)-powers are asymptotic scales, not a large--sieve heuristic.

\subsection{Fixed target orthogonality does not survive target synthesis}

For one fixed \(N\), the functions
\begin{equation}
 e_{\chi,N}(F)=\chi(N+F)\overline{\chi(N)}
 \label{eq:tpc38-intro-echi}
\end{equation}
are orthogonal over \(F\bmod M\).  This recovers the one--period masses
\begin{equation}
 \sum_{F\bmod M}|C_N(F)|^2=A,\qquad
 \sum_{F\bmod M}|P_N(F)|^2=1-A,
 \label{eq:tpc38-intro-fixed-target-masses}
\end{equation}
where
\begin{equation}
 A=\prod_{i=1}^3\left(1-\frac1{q_i-1}\right).
 \label{eq:tpc38-intro-A}
\end{equation}
When \(N\) moves, the cross Gram becomes a product of shifted Jacobi
sums.  More directly, the full face rows satisfy
\begin{equation}
 \langle C_N,C_{N'}\rangle
 =\prod_{i=1}^3
 \begin{cases}
  1-p_i,&N\equiv N'\pmod{q_i},\\
  1-p_i-p_i^2,&N\not\equiv N'\pmod{q_i},
 \end{cases}
 \label{eq:tpc38-intro-full-row-Gram}
\end{equation}
with \(p_i=(q_i-1)^{-1}\).  Distinct target rows therefore have inner
product \(1-O(J^{-1})\) when they are distinct in all coordinates.  They
are almost parallel, not almost orthogonal.

We compute the entire scalar target--synthesis spectrum.  For
\begin{equation}
 (T_Ca)(F)=\sum_{N\bmod M}^{*}a_NC_N(F),
 \label{eq:tpc38-intro-TC}
\end{equation}
the squared singular values are indexed by target Mellin modes.  The
largest is
\begin{equation}
 \|T_C\|_{\op}^2
 =\prod_{i=1}^3\frac{q_i(q_i-2)^2}{(q_i-1)^2}
 \asymp\Phi\asymp J^3.
 \label{eq:tpc38-intro-TC-norm}
\end{equation}
For the proper union, the principal squared singular value is
\(\asymp\Phi/J\asymp J^2\); all its nonprincipal target modes are
\(O(1)\).  The small fixed--row norm \(1-A\asymp J^{-1}\) therefore does
not imply a small synthesis norm.

\subsection{Recombination and the identity bucket}

The full and proper character families are complementary.  Summing them
before estimating gives
\begin{equation}
 \boxed{C_N(F)+P_N(F)=\one_{M\mid F}.}
 \label{eq:tpc38-intro-recombination}
\end{equation}
Every nonjoint ghost cancels.  Separate estimates for \(C_N\) and \(P_N\)
are consequently sufficient but strictly stronger than an estimate for
their recombined congruence completion.

In the physical determinant interval, \eqref{eq:tpc38-intro-recombination}
leaves only
\begin{equation}
 F=kM,\qquad |k|\ll X/M.
 \label{eq:tpc38-intro-joint-aliases}
\end{equation}
At the endpoint
\begin{equation}
 Q=X^{267/400+o(1)},\qquad
 J=X^{133/400+o(1)},\qquad
 M=X^{399/400+o(1)},
 \label{eq:tpc38-intro-endpoint}
\end{equation}
the endpoint alias order is
\begin{equation}
 \cK_X:=1+\frac XM
 =\frac Q{J^2}X^{o(1)}
 =X^{1/400+o(1)}.
 \label{eq:tpc38-intro-KX}
\end{equation}
On \(F=kM\), the ratio in \eqref{eq:tpc38-intro-ratio} is \(1\) and every
character phase in \eqref{eq:tpc38-intro-character-face} equals one.
Multiplicative character orthogonality is therefore exactly blind on the
remaining bucket.

Write \(\mathscr Z_k\) for the raw shifted output on
\(su=N+kM\).  The recombined completed output is
\begin{equation}
 \left\|\sum_k\mathscr Z_k\right\|_2^2
 =\sum_{k,k'}\langle\mathscr Z_k,\mathscr Z_{k'}\rangle.
 \label{eq:tpc38-intro-alias-Gram}
\end{equation}
If its alias columns satisfy the diagonal trace bound
\begin{equation}
 \sum_k\|\mathscr Z_k\|_2^2
 \ll X^{o(1)}Q^2J\cK_X,
 \label{eq:tpc38-intro-diagonal-input}
\end{equation}
Cauchy gives only \(X^{o(1)}Q^2J\cK_X^2\).  The desired scale is
\begin{equation}
 Q^2J\cK_X=\frac{Q^3}{J}X^{o(1)}.
 \label{eq:tpc38-intro-target-scale}
\end{equation}
The exact finite loss is the alias count \(K_B\); at the endpoint it is of
order \(\cK_X\), and a rank--one family of alias columns attains it.

There is a more basic logical distinction.  The original physical output
is \(\mathscr Z_0\), not the alias sum in
\eqref{eq:tpc38-intro-alias-Gram}.  No norm bound for the latter recovers
the former: for two aliases, the abstract choice
\(\mathscr Z_0=v\), \(\mathscr Z_1=-v\) makes the alias sum and both
fixed full/proper field outputs vanish while \(\mathscr Z_0\) is arbitrary.
Thus the shifted--fiber estimate proposed in TPC--37 controls the
congruence--completed output but needs an additional equality--recovery
input before it can imply the original orbit gate.  A uniform column bound
at scale \(Q^2J\) would recover \(\mathscr Z_0\), but at \(k=0\) it is
already stronger than the desired \(Q^2J\cK_X\) bound.  A full alias--Gram
bound is stronger still and automatically contains this column bound.

\subsection{Arithmetic meaning and claim boundary}

On the terminal submodel \(s=1\), a joint alias contains
\begin{equation}
 a_N(k)=-\mu(N+kM)\log(N+kM)\Phi_0(N+kM).
 \label{eq:tpc38-intro-terminal-sequence}
\end{equation}
Its squared alias sum has cross terms
\begin{equation}
 \mu(N+kM)\mu(N+(k+r)M),
 \label{eq:tpc38-intro-Mobius-correlation}
\end{equation}
and restoring the physical rows gives the affine four--M\"obius Gram
identified in TPC--34.  This is an averaged, highly structured
correlation.  It is not logically equivalent to pointwise Chowla in each
short progression.

There is also a necessary distinction in terminology.  Bounds for the two
prescribed fields \(C\) and \(P\) do not constitute an all--field operator
estimate.  An estimate of the same normalization for every coefficient
field, stable under fiber restriction, is strictly stronger.  Since
\(C_N(0)=A=1-O(J^{-1})\), that hereditary all--field estimate
restricted to \(F=0,s=1\) would imply a diagonal--scale estimate for the
original terminal operator, stronger than the inherited target by the
factor \(\cK_X\).  We prove no hereditary all--field theorem, positive
physical equality--recovery estimate, or M\"obius correlation estimate.

The unconditional conclusions of the paper are the exact character
projectors, shifted Jacobi Grams, ratio Parseval theorem, exact collision
multiplicity and its scale ladder,
target--synthesis spectrum, face recombination, equality--nonrecoverability,
and sharp abstract coherence examples.  They locate the remaining gates
but do not prove the physical orbit estimate, affine Chowla, sieve--parity failure, a
Hardy--Littlewood asymptotic, or infinitely many twin primes.

\Cref{sec:tpc38-interface} fixes the inherited physical normalization.
\Cref{sec:tpc38-character-projectors} proves the target--coupled character
calculus.  \Cref{sec:tpc38-ratio-parseval} gives the ratio and determinant
stratifications.  \Cref{sec:tpc38-target-spectrum} computes the synthesis
spectrum.  \Cref{sec:tpc38-recombined-gate} reduces to the identity bucket.
\Cref{sec:tpc38-terminal-boundary} records its arithmetic content, and
\cref{sec:tpc38-next-gate,sec:tpc38-certificate-claims} state the next
inputs and formal claim boundary.
